% -----------------------------------------------------------
\section{Godliness, Godly, Ungodly}
% -----------------------------------------------------------
	\label{GODLINESS}
	\label{UNGODLY}
	
	See also \hyperref[MYSTERY]{Mystery}, 
	
% % ----------------------
% \subsection{See also}
% % ----------------------

% ----------------------
\subsection{1828 American Dictionary of the English Language} % *1828
% ----------------------
	\label{GODLINESS-1828}
	
	% -----
	\subsubsection{Godliness}
	% -----

		\begin{compactitem}
			\item Piety; belief in God, and reverence for his character and laws.
			\item[1] A religious life; a careful observance of the laws of God and performance of religious duties, proceeding from love and reverence for the divine character and commands; christian obedience.
			\item[2] Revelation; the system of christianity.
		\end{compactitem}
		
	% -----
	\subsubsection{Godly}
	% -----

		\begin{compactitem}
			\item Pious; reverencing God, and his character and laws.
			\item[1] Living in obedience to God's commands, from a principle of love to him and reverence of his character and precepts; religious; righteous; as a \textit{godly} person.
			\item[2] Pious; conformed to God's law; as a \textit{godly} life.
		\end{compactitem}

% ----------------------
\subsection{Oxford English Dictionary} % *OED
% ----------------------
	\label{GODLINESS-OED}
	
	% -----
	\subsubsection{Godliness}
	% -----

		\begin{compactitem}
			\item[1] The quality of being godly; devout observance of the law of God; piety.
			\item[2] With possessive adjective (as \textit{\textbf{your Godliness}}), used as a title of respect.
		\end{compactitem}
		
	% -----
	\subsubsection{Godly (adjective)}
	% -----

		\begin{compactitem}
			\item[1] Of or relating to God; coming from God; divine; spiritual. Now somewhat \textit{archaic}.
			\item[2] Devoutly observant of the laws of God; religious, pious.
			\item[2.a] Of conduct, speech, etc. Now somewhat \textit{archaic}.
			\item[2.b] Of a person.
			\item[2.b.(a)] In general use.
			\item[2.b.(b)] Designating puritan supporters, typically Presbyterians and Independents (Congregationalists), of the parliamentary cause in the English Civil War, who gained political ascendancy during the period of the Commonwealth and the Cromwellian Protectorate with the aim of spiritually and morally reforming the nation; (later, by extension) designating zealous Nonconfomists or Low Church Anglicans who followed similar principles and practices. Now \textit{historical}.
		\end{compactitem}

% ----------------------
\subsection{Scriptural Usage} % *SCRIPTURES
% ----------------------
	\label{GODLINESS-SCRIPTURES}

	% -----
	\subsubsection{Old Testament} % *OT
	% -----
		\label{GODLINESS-OT}
	
		% --
		\paragraph{KJV}
		% --
			\label{GODLINESS-OT-KJV}
	
			\begin{tabular}{ll}
				Total occurrences in the OT & 4
			\end{tabular}
			
			\begin{compactdesc}
				\item[{\hyperref[PSALMS-4-3]{Psalms 4:3}}]
				\item[{\hyperref[PSALMS-12-1]{Psalms 12:1}}]
				\item[{\hyperref[PSALMS-32-6]{Psalms 32:6}}]
				\item[{\hyperref[MALACHI-2-15]{Malachi 2:15}}]
			\end{compactdesc}
			
		% % --
		% \paragraph{Hebrew Bible}
		% % --
			% \label{GODLINESS-HEBREW}
		
			% %
			% \subparagraph{\texorpdfstring{\he{sample word}}{}}
			% %
				% \label{PAGA} % whatever is in step bible without periods
			
				% \begin{compactdesc}
					% \item[HALOT , PDF ] \textbf{}
						% \begin{compactitem}
							% \item[1]
						% \end{compactitem}
					% \item[BDB , PDF ] \textbf{}
						% \begin{compactitem}
							% \item[1]
						% \end{compactitem}
					% \item[DCH PDF ] \textbf{}
						% \begin{compactitem}
							% \item[1]
						% \end{compactitem}
				% \end{compactdesc}
				
				% \begin{compactdesc}
					% \item[Some OT uses] \textbf{}
						% \begin{compactdesc}
							% \item[]
						% \end{compactdesc}
				% \end{compactdesc}
			
		% % --
		% \paragraph{Septuagint}
		% % --
			% \label{GODLINESS-LXX}
		
			% %
			% \subparagraph{\texorpdfstring{\gr{sample word}}{}}
			% %
		
	% -----
	\subsubsection{New Testament} % *NT
	% -----
		\label{GODLINESS-NT}
	
		% --
		\paragraph{KJV}
		% --
			\label{GODLINESS-NT-KJV}
			
	
			\begin{tabular}{ll}
				Total occurrences in the NT & 26
			\end{tabular}
			
			\begin{compactdesc}
				\item[{\hyperref[2CORINTHIANS-1-12]{2 Corinthians 1:12}}] 
				\item[{\hyperref[2CORINTHIANS-7-9]{2 Corinthians 7:9}}]
				\item[{\hyperref[2CORINTHIANS-7-10]{2 Corinthians 7:10}}]
				\item[{\hyperref[2CORINTHIANS-7-11]{2 Corinthians 7:11}}]
				\item[{\hyperref[2CORINTHIANS-11-2]{2 Corinthians 11:2}}]
				\item[{\hyperref[1TIMOTHY-1-4]{1 Timothy 1:4}}]
				\item[{\hyperref[1TIMOTHY-2-2]{1 Timothy 2:2}}]
				\item[{\hyperref[1TIMOTHY-2-10]{1 Timothy 2:10}}]
				\item[{\hyperref[1TIMOTHY-3-16]{1 Timothy 3:16}}]
				\item[{\hyperref[1TIMOTHY-4-7]{1 Timothy 4:7}}]
				\item[{\hyperref[1TIMOTHY-4-8]{1 Timothy 4:8}}]
				\item[{\hyperref[1TIMOTHY-6-3]{1 Timothy 6:3}}]
				\item[{\hyperref[1TIMOTHY-6-5]{1 Timothy 6:5}}]
				\item[{\hyperref[1TIMOTHY-6-6]{1 Timothy 6:6}}]
				\item[{\hyperref[1TIMOTHY-6-11]{1 Timothy 6:11}}]
				\item[{\hyperref[2TIMOTHY-3-5]{2 Timothy 3:5}}]
				\item[{\hyperref[2TIMOTHY-3-12]{2 Timothy 3:12}}]
				\item[{\hyperref[TITUS-1-1]{Titus 1:1}}]
				\item[{\hyperref[TITUS-2-12]{Titus 2:12}}]
				\item[{\hyperref[HEBREWS-12-28]{Hebrews 12:28}}]
				\item[{\hyperref[2PETER-1-3]{2 Peter 1:3}}]
				\item[{\hyperref[2PETER-1-6]{2 Peter 1:6}}]
				\item[{\hyperref[2PETER-1-7]{2 Peter 1:7}}]
				\item[{\hyperref[2PETER-2-9]{2 Peter 2:9}}]
				\item[{\hyperref[2PETER-3-11]{2 Peter 3:11}}]
				\item[{\hyperref[3JOHN-1-6]{3 John 1:6}}]
			\end{compactdesc}
			
		% --
		\paragraph{Greek New Testament} % *GREEK
		% --
			\label{GODLINESS-GREEK}
			
			%
			\subparagraph{\texorpdfstring{\gr{ἀσεβής}}{}}
			%
				\label{ASEBES}
			
				\begin{compactdesc}
					\item[BDAG 123b] \textbf{}
						\begin{compactitem}
							\item[1] \textbf{pertaining to violating norms for a proper relation to deity, \textit{irreverent, impious, ungodly}}
							\item[1.a] of persons
							\item[1.b] of human characteristics
						\end{compactitem}
					% \item[CGL , PDF ] \textbf{}
						% \begin{compactitem}
							% \item 
						% \end{compactitem}
					% \item[LSJ , PDF ] \textbf{}
						% \begin{compactitem}
							% \item 
						% \end{compactitem}
				\end{compactdesc}
				
				% \begin{tabular}{ll}
					% Total occurrences in the NT & 
				% \end{tabular}
				
				% \begin{compactdesc}
					% \item[Some NT uses] \textbf{}
						% \begin{compactdesc}
							% \item[]
						% \end{compactdesc}
				% \end{compactdesc}
				
			% %
			% \subparagraph{TDNT: \texorpdfstring{\gr{sample word}}{}  (3:536--556)}
			% %
				% \label{KATALLAGE-TDNT}
		
			%
			\subparagraph{\texorpdfstring{\gr{εὐσέβεια}}{}}
			%
				\label{EUSEBEIA}
			
				\begin{compactdesc}
					\item[BDAG 363b] \textbf{}
						\begin{compactitem}
							\item in our literature and in the LXX only of \textbf{awesome respect accorded to God, \textit{devoutness, piety, godliness}}
						\end{compactitem}
					\item[CGL 630a, PDF 656] \textbf{}
						\begin{compactitem}
							\item[1] reverent conduct (towards gods or parents), \textbf{reverence, piety}
							\item[2] (as a deity) \textbf{Piety}
						\end{compactitem}
					\item[LSJ , PDF ] \textbf{}
						\begin{compactitem}
							\item 
						\end{compactitem}
				\end{compactdesc}
				
				\begin{compactdesc}
					\item[Some NT uses] \textbf{}
						\begin{compactdesc}
							\item[{\hyperref[1TIMOTHY-2-2]{1 Timothy 2:2}}]
							\item[{\hyperref[1TIMOTHY-3-16]{1 Timothy 3:16}}]
							\item[{\hyperref[1TIMOTHY-4-7]{1 Timothy 4:7}}]
							\item[{\hyperref[1TIMOTHY-4-8]{1 Timothy 4:8}}]
							\item[{\hyperref[1TIMOTHY-6-3]{1 Timothy 6:3}}]
							\item[{\hyperref[1TIMOTHY-6-5]{1 Timothy 6:5}}]
							\item[{\hyperref[1TIMOTHY-6-6]{1 Timothy 6:6}}]
							\item[{\hyperref[1TIMOTHY-6-11]{1 Timothy 6:11}}]
							\item[{\hyperref[2TIMOTHY-3-5]{2 Timothy 3:5}}]
							\item[{\hyperref[TITUS-1-1]{Titus 1:1}}]
							\item[{\hyperref[2PETER-1-3]{2 Peter 1:3}}]
							\item[{\hyperref[2PETER-1-6]{2 Peter 1:6}}]
							\item[{\hyperref[2PETER-1-7]{2 Peter 1:7}}]
							\item[{\hyperref[2PETER-3-11]{2 Peter 3:11}}]
						\end{compactdesc}
				\end{compactdesc}
				
			%
			\subparagraph{\texorpdfstring{\gr{θεοσέβεια}}{}}
			%
				\label{THEOSEBEIA}
				
				In the NT, the only use of this word is at \hyperref[1TIMOTHY-2-10]{1 Timothy 2:10}.
			
				\begin{compactdesc}
					\item[BDAG 400a] \textbf{}
						\begin{compactitem}
							\item reverence for God or set of beliefs and practices relating to interest in God, \textit{piety, godliness}
						\end{compactitem}
					\item[CGL 679b, PDF 705] \textbf{}
						\begin{compactitem}
							\item reverence for the gods, \textbf{piety}
						\end{compactitem}
					\item[LSJ , PDF ] \textbf{}
						\begin{compactitem}
							\item 
						\end{compactitem}
				\end{compactdesc}
				
				\begin{compactdesc}
					\item[Some NT uses] \textbf{}
						\begin{compactdesc}
							\item[{\hyperref[1TIMOTHY-2-10]{1 Timothy 2:10}}]
						\end{compactdesc}
				\end{compactdesc}
				
			%
			\subparagraph{\texorpdfstring{\gr{θειότης}}{}}
			%
				\label{THEIOTES}
			
				\begin{compactdesc}
					\item[BDAG 395b] (of a divinity: the term in such description is not tautologous but usually refers to performance that one might properly associate with a divinity)
						\begin{compactitem}
							\item \textbf{the quality or characteristic(s) pertaining to deity, \textit{divinity, divine nature, divineness}}
						\end{compactitem}
					% \item[CGL , PDF ] \textbf{}
						% \begin{compactitem}
							% \item 
						% \end{compactitem}
					% \item[LSJ , PDF ] \textbf{}
						% \begin{compactitem}
							% \item 
						% \end{compactitem}
				\end{compactdesc}
				
				\begin{tabular}{ll}
					Total occurrences in the NT & 1
				\end{tabular}
				
				\begin{compactdesc}
					\item[Some NT uses] \textbf{}
						\begin{compactdesc}
							\item[{\hyperref[ROMANS-1-20]{Romans 1:20}}]
						\end{compactdesc}
				\end{compactdesc}
				
			%
			\subparagraph{\texorpdfstring{\gr{θεότης}}{}}
			%
				\label{THEOTES}
			
				\begin{compactdesc}
					\item[BDAG 400b] \textbf{}
						\begin{compactitem}
							\item \textbf{the state of being god, \textit{divine character/nature, deity, divinity}}, used as abstract noun for \gr{θεός}
						\end{compactitem}
				\end{compactdesc}
		
	% -----
	\subsubsection{Book of Mormon} % *BOM
	% -----
		\label{GODLINESS-BOM}
	
		\begin{tabular}{ll}
			Total occurrences in the BOM & 1
		\end{tabular}
		
		\begin{compactdesc}
			\item[{\hyperref[MORONI-7-30]{Moroni 7:30}}]
		\end{compactdesc}
		
	% -----
	\subsubsection{Doctrine and Covenants} % *DC
	% -----
		\label{GODLINESS-DC}
	
		\begin{tabular}{ll}
			Total occurrences in the DC & 6
		\end{tabular}
		
		\begin{compactdesc}
			\item[{\hyperref[DC4-6]{DC 4:6}}]
			\item[{\hyperref[DC19-10]{DC 19:10}}]
			\item[{\hyperref[DC20-69]{DC 20:69}}]
			\item[{\hyperref[DC84-20]{DC 84:20}}] 
			\item[{\hyperref[DC84-21]{DC 84:21}}] 
			\item[{\hyperref[DC107-30]{DC 107:30}}]
		\end{compactdesc}
		
	% -----
	\subsubsection{Pearl of Great Price} % *PGP
	% -----
		\label{GODLINESS-PGP}
	
		\begin{tabular}{ll}
			Total occurrences in the PGP & 1
		\end{tabular}
		
		\begin{compactdesc}
			\item[JS-H 1:19]
		\end{compactdesc}
		
% ----------------------
\subsection{Key Phrases} % *KEYPHRASES *PHRASES
% ----------------------
	\label{GODLINESS-PHRASES}

	\begin{compactdesc}
	
		\item[mystery of godliness] \textbf{2} | \hyperref[1TIMOTHY-3-16]{1 Timothy 3:16}; \hyperref[DC19-10]{DC 19:10}
			\begin{compactdesc}
				\item[Bible anchor verses] \phantom{.}
					\begin{compactdesc}
						\item[1 Timothy 3:16] \gr{καὶ ὁμολογουμένως μέγα ἐστὶν τὸ τῆς εὐσεβείας μυστήριον·  Ὃς ἐφανερώθη ἐν σαρκί, ἐδικαιώθη ἐν πνεύματι, ὤφθη ἀγγέλοις, ἐκηρύχθη ἐν ἔθνεσιν, ἐπιστεύθη ἐν κόσμῳ, ἀνελήμφθη ἐν δόξῃ.}
					\end{compactdesc}
				\item[Commentary] See \hyperref[GODLINESS-mystery]{this study} below.
			\end{compactdesc}
			
		\item[power of godliness] 
		
	\end{compactdesc}
	
% % ----------------------
% \subsection{Bible Dictionaries} % *DICTIONARIES
% % ----------------------
	% \label{GODLINESS-BD}

% ----------------------
\subsection{Restoration Usage} % *RESTORATION
% ----------------------
	\label{GODLINESS-RESTORATION}

	% -----
	\subsubsection{Joseph Smith} % *JS
	% -----
		\label{GODLINESS-JS}
	
		\begin{compactdesc}
			\item[{\hyperref[EMS-1842-JS-9-JUN-1842]{9 June 1842}}] It is one evidence that men are unacquainted with the principle of godliness, to behold the contraction of feeling and lack of charity. The pow'r and glory of Godliness is spread out on a broad principle to throw out the mantle of charity. God does not look on sin with allowance, but when men have sin'd there must be allowance made for them.
		\end{compactdesc}
		
	% -----
	\subsubsection{Temple Ordinances}
	% -----
		\label{GODLINESS-TEMPLE}
	
		\begin{compactdesc}
			\item[{\hyperref[TEMPLE-ENDOWMENT-ENDOW-PRE-1990-LECTURE]{Newer lecture at the veil}}] These are what are termed the mysteries of godliness---that which will enable you to understand the expression of the Savior, made just prior to his betrayal: ``This is life eternal, that they might know thee, the only true God, and Jesus Christ, whom thou has sent.''
		\end{compactdesc}
	
	% % -----
	% \subsubsection{Brigham Young} % *BY
	% % -----
		% \label{GODLINESS-BY}
	
		% \begin{compactdesc}
			% \item[{\hyperref[KEY]{Date/Source}}]
		% \end{compactdesc}
	
% % ----------------------
% \subsection{Commentary and Studies} % *COMMENTARY *STUDIES
% % ----------------------
	% \label{GODLINESS-COMMENTARY}

	% % -----
	% \subsubsection{Key Points}
	% % -----
		% \label{GODLINESS-KEYS}
	
		% \begin{compactitem}
			% \item
		% \end{compactitem}
		
	% -----
	\subsubsection{The ``mystery of godliness''}
	% -----
		\label{GODLINESS-mystery}